\documentclass[10pt,a4paper]{amsart}
\usepackage[margin=19mm]{geometry}
\usepackage{amsmath,amssymb,mathtools}
\usepackage[scr=rsfs,cal=euler]{mathalfa}
\usepackage{xfrac}
\usepackage[unicode,hidelinks,bookmarks=false]{hyperref}
\newcommand{\ie}{i.e.\ }
\newcommand{\cf}{cf.\ }
\newcommand{\resp}{respectively\ }
\newcommand{\Subsection}{\subsection}
\providecommand{\nobreakdash}{\nobreak\hbox{-}}
\setlength{\emergencystretch}{2em}
\multlinegap0pt
\usepackage{bm}
\usepackage{bbm}
\usepackage{dsfont}
\usepackage{xspace}
\usepackage{cancel}
\usepackage{mathtools}
\usepackage{amsthm}
\makeatletter
\@ifundefined{theorem}{\newtheorem{theorem}{Theorem}}{}
\@ifundefined{assumption}{\newtheorem{assumption}[theorem]{Assumption}}{}
\@ifundefined{lemma}{\newtheorem{lemma}[theorem]{Lemma}}{}
\makeatother
\def\eqref#1{equation~\ref{#1}}
\title[A Decision-Theoretic View of Test-Time Training: When, How F]{A Decision-Theoretic View of Test-Time Training: When, How Far, and Which Directions to Adapt}
\author{Tomoya Wakayama}
\date{Source: arXiv:2606.15569v1 (CC BY 4.0)}
\begin{document}
\maketitle

\noindent\emph{Source and licence.} A Decision-Theoretic View of Test-Time Training: When, How Far, and Which Directions to Adapt. Version arXiv:2606.15569v1 (CC BY 4.0), distributed under the Creative Commons Attribution 4.0 International licence, as recorded in the arXiv licence metadata for this exact version and shown on the abstract page https://arxiv.org/abs/2606.15569v1. Full rights notice: licenses/arxiv-2606.15569-CC-BY-4.0.txt.\par\smallskip

\noindent\textit{Editorial scope.} Counted unit: the Theorem whose source label is \texttt{thm:finite\_step\_deviation\_prompt\_app} in the article (source0.tex lines 2852--2865, source label \texttt{thm:finite\_step\_deviation\_prompt\_app}), delivered together with the article's own complete original proof of it (source0.tex lines 2867--2950, one \texttt{proof} environment of 84 lines). No mathematics is written, repaired or re-derived: statement and proof are the article's own text line for line. Required context retained uncounted: eq:app\_true\_gd (lines 2761--2766); eq:app\_lin\_gd (lines 2772--2778); eq:app\_lin\_kernel\_recursion (lines 2781--2785); assump:JL\_local\_app (lines 2796--2806); eq:app\_J\_Lipschitz (lines 2800--2805); eq:app\_J\_bounded (lines 2800--2805); eq:app\_rho\_small (lines 2822--2825); assump:iterate\_stability\_app (lines 2829--2833); lem:nonlinear\_prompt\_recursion\_app (lines 2837--2840). Every definition, display and label of those lines is retained unchanged. Omitted from this package and not redistributed: 1--2760, 2767--2771, 2779--2780, 2786--2795, 2806--2821, 2826--2828, 2834--2836, 2841--2851, 2866--2866, 2951--3907. Nothing inside a retained excerpt is omitted or altered: every retained body is delivered exactly as the article's lines are written, so no omission receipt of any kind accompanies this package. Citations: the retained excerpts carry no citation command at all, so no bibliography entry of the article is transcribed and no reference is redistributed. Reference adaptation: none was needed and none was made; every reference inside the delivered unit resolves inside it, so no unresolved reference or citation is produced. Printed slips kept literal: no printed word, symbol, hypothesis or number of the counted statement or of its proof was altered. Layout and class: the article's own class is \texttt{article} with packages including microtype, graphicx, subcaption, booktabs, natbib, enumitem, tikz, amsmath, amsfonts, bm, hyperref, needspace, mathtools, icml2026; the wrapper is the lane's pinned \texttt{amsart} editorial wrapper at 10pt on A4, which restates the theorem-like environments the retained text needs and adds the article's own macro definitions that the retained text needs (\texttt{\textbackslash eqref}), with extra wrapper packages (bm, bbm, dsfont, xspace, cancel, mathtools). The delivered numbering is the unit's own. Assets: no retained span contains a figure environment, a graphics-inclusion command, a PDF-page inclusion, a table or a TikZ picture; no image file is copied into this package, the package declares no asset dependency and no image file is redistributed with it. Licence: version arXiv:2606.15569v1 of this article is distributed under the Creative Commons Attribution 4.0 International licence, as recorded in the arXiv licence metadata for this exact version and shown on the abstract page https://arxiv.org/abs/2606.15569v1. A full rights notice is delivered at licenses/arxiv-2606.15569-CC-BY-4.0.txt. Characters: every retained excerpt is pure ASCII, so the delivered engine, XeLaTeX with its default Latin Modern text font, reports no missing character.
\begin{equation}
    \theta_{t+1} 
    =\theta_t-\eta\nabla L(\theta_t) 
    =\theta_t+\rho\,J(\theta_t)^\top\bigl(r-s(\theta_t)\bigr),
\label{eq:app_true_gd}
\end{equation}

\begin{equation}
    \theta_{t+1}^{\mathrm{lin}} 
    =\theta_t^{\mathrm{lin}}+\rho\,\Phi^\top\bigl(r-\Phi\theta_t^{\mathrm{lin}}\bigr), 
    \quad
    \theta_0^{\mathrm{lin}}=0,
    \label{eq:app_lin_gd}
\end{equation}

\begin{equation}
f_{t+1}^{\mathrm{lin}}=f_t^{\mathrm{lin}}+\rho K_0\,(r-f_t^{\mathrm{lin}})
=f_t^{\mathrm{lin}}+\rho K_0\,e_t^{\mathrm{lin}},
\label{eq:app_lin_kernel_recursion}
\end{equation}

\begin{assumption}[Bounded and Lipschitz Jacobian]
\label{assump:JL_local_app}
There exist $R>0$ and constants $B,L>0$ such that for all $\theta,\theta'\in\mathbb{R}^d$ with
$\|\theta\|_2\le R$ and $\|\theta'\|_2\le R$,
\begin{align}
\|J(\theta)\|_{\mathrm{op}} &\le B,
\label{eq:app_J_bounded}\\
\|J(\theta)-J(\theta')\|_{\mathrm{op}} &\le L\,\|\theta-\theta'\|_2.
\label{eq:app_J_Lipschitz}
\end{align}
\end{assumption}

\begin{equation}
    0<\rho\le \frac{1}{B^2},
\label{eq:app_rho_small}
\end{equation}

\begin{assumption}[Iterate stability]
\label{assump:iterate_stability_app}
For a given $T\in\mathbb{N}$, the nonlinear GD iterates \eqref{eq:app_true_gd} satisfy
$\|\theta_t\|_2\le R$ for all $t=0,1,\dots,T$.
\end{assumption}

\begin{lemma}[Nonlinear prompt-space recursion]
\label{lem:nonlinear_prompt_recursion_app}
Under Assumption~\ref{assump:JL_local_app}, if $\|\theta_t\|_2\le R$ and $\|\theta_{t+1}\|_2\le R$, then $f_{t+1}=f_t+\rho K_t e_t + \xi_t$, where the remainder term $\xi_t\in\mathbb{R}^n$ satisfies $\|\xi_t\|_2 \ \le\ \frac{L}{2}\,\|\theta_{t+1}-\theta_t\|_2^2$.
\end{lemma}

\begin{theorem}[Finite-step deviation between nonlinear and linearized prompt corrections]\label{thm:finite_step_deviation_prompt_app}
Fix a $T\in\mathbb{N}$. Assume Assumptions~\ref{assump:JL_local_app} and \ref{assump:iterate_stability_app},
and the step-size condition \eqref{eq:app_rho_small}. Let $f_T=s(\theta_T)$ be the nonlinear prompt correction after
$T$ steps of \eqref{eq:app_true_gd}, and let $f_T^{\mathrm{lin}}=\Phi\theta_T^{\mathrm{lin}}$ be the linearized prompt
correction after $T$ steps of \eqref{eq:app_lin_gd}. Define the residual scale inside the stability ball $E_R:=\sup_{\|\theta\|_2\le R}\|r-s(\theta)\|_2.$
Then
\begin{align}
\|f_T-f_T^{\mathrm{lin}}\|_2 &\le \frac{3}{2}\,L\,B^2\,\rho^2\,T^2\, \bigl(\max\{E_R,\|r\|_2\}\bigr)^2 \\
&\le \frac{3}{2}\,L\,B^2\,\rho^2\,T^2\,E_R^2.
\label{eq:app_finite_step_deviation_bound}
\end{align}
In particular, for fixed local regularity constants $(B,L)$ and residual scale $E_R$, the discrepancy is
$O\bigl((\rho T)^2\bigr)$.
\end{theorem}

\begin{proof}
Define the discrepancy $\delta_t:=f_t-f_t^{\mathrm{lin}}\in\mathbb{R}^n$. By
Lem.~\ref{lem:nonlinear_prompt_recursion_app} and \eqref{eq:app_lin_kernel_recursion},
\begin{align*}
\delta_{t+1}
&=
\bigl(f_t+\rho K_t e_t+\xi_t\bigr)-\bigl(f_t^{\mathrm{lin}}+\rho K_0 e_t^{\mathrm{lin}}\bigr) \\
&=
\delta_t+\rho K_t(r-f_t)-\rho K_0(r-f_t^{\mathrm{lin}})+\xi_t \\
&=
\delta_t+\rho K_t(e_t^{\mathrm{lin}}-\delta_t)-\rho K_0 e_t^{\mathrm{lin}}+\xi_t \\
&=
(I_n-\rho K_t)\delta_t + \rho (K_t-K_0)e_t^{\mathrm{lin}}+\xi_t.
\end{align*}
Taking norms and using $\|I_n-\rho K_t\|_{\mathrm{op}}\le 1$ (from \eqref{eq:app_rho_small} and $\|J(\theta_t)\|_{\mathrm{op}}\le B$),
\begin{equation}
\|\delta_{t+1}\|_2
\le
\|\delta_t\|_2 + \rho\|K_t-K_0\|_{\mathrm{op}}\|e_t^{\mathrm{lin}}\|_2+\|\xi_t\|_2.
\label{eq:app_delta_step}
\end{equation}

Under \eqref{eq:app_rho_small}, $0\preceq I_n-\rho K_0\preceq I_n$, so
$e_t^{\mathrm{lin}}=(I_n-\rho K_0)^t r$ implies $\|e_t^{\mathrm{lin}}\|_2\le \|r\|_2$ for all $t$.

Write $J_t:=J(\theta_t)$. Then, $K_t-K_0=J_tJ_t^\top-J_0J_0^\top=(J_t-J_0)J_t^\top + J_0(J_t-J_0)^\top,$
so
\(
\|K_t-K_0\|_{\mathrm{op}}
\le \|J_t-J_0\|_{\mathrm{op}}\|J_t\|_{\mathrm{op}}+\|J_0\|_{\mathrm{op}}\|J_t-J_0\|_{\mathrm{op}}
\le 2B\,\|J_t-J_0\|_{\mathrm{op}}.
\)
By \eqref{eq:app_J_Lipschitz}, $\|J_t-J_0\|_{\mathrm{op}}\le L\|\theta_t\|_2$. Moreover, from \eqref{eq:app_true_gd} and
\eqref{eq:app_J_bounded}, we have
\begin{align*}
   \|\theta_{t+1}-\theta_t\|_2 
   &=\rho\|J(\theta_t)^\top e_t\|_2  \\ 
   &\le \rho\,\|J(\theta_t)\|_{\mathrm{op}}\,\|e_t\|_2 \\
   &\le \rho\,B\,E_R,
\end{align*}
since $\|e_t\|_2=\|r-s(\theta_t)\|_2\le E_R$ for $\|\theta_t\|_2\le R$.
Thus $\|\theta_t\|_2\le t\,\rho B E_R$, and hence
\begin{equation}
    \|K_t-K_0\|_{\mathrm{op}} \le 2B  L\|\theta_t\|_2 \le 2LB^2\,\rho\,t\,E_R.
\label{eq:app_K_drift}
\end{equation}


From Lem.~\ref{lem:nonlinear_prompt_recursion_app} and the bound on $\|\theta_{t+1}-\theta_t\|_2$,
\begin{equation}
     \|\xi_t\|_2
     \le \frac{L}{2}\,(\rho B E_R)^2
     =\frac{L}{2}\,B^2\rho^2 E_R^2.
\label{eq:app_xi_final}
\end{equation}

Plug \eqref{eq:app_K_drift}, $\|e_t^{\mathrm{lin}}\|_2\le \|r\|_2$, and \eqref{eq:app_xi_final} into \eqref{eq:app_delta_step}:
\[
\|\delta_{t+1}\|_2
\le
\|\delta_t\|_2
+
\rho  (2LB^2\rho t E_R)  \|r\|_2
+
\frac{L}{2}B^2\rho^2 E_R^2.
\]
Let $E:=\max\{E_R,\|r\|_2\}$. Then $\|r\|_2\le E$ and $E_R\le E$, so
\[
\|\delta_{t+1}\|_2
\le
\|\delta_t\|_2
+
2LB^2\rho^2\,t\,E^2
+
\frac{L}{2}B^2\rho^2 E^2.
\]
Since $\delta_0=0$, summing from $t=0$ to $T-1$ gives
\begin{align*}
    \|\delta_T\|_2 &\le 2LB^2\rho^2 E^2 \sum_{t=0}^{T-1} t + \frac{L}{2}B^2\rho^2 E^2 \sum_{t=0}^{T-1} 1 \\
    &= LB^2\rho^2 E^2 T(T-1) + \frac{L}{2}B^2\rho^2 E^2 T \\
    &\le \frac{3}{2}LB^2\rho^2 E^2 T^2,
\end{align*}
which is \eqref{eq:app_finite_step_deviation_bound}.
\end{proof}

\end{document}
